\documentclass[12pt]{article}
\usepackage{amsmath,amssymb,amsthm}
\usepackage[margin=1in]{geometry}

\title{Problem 118: Pandigital Prime Sets}
\author{Project Euler Solution}
\date{}

\begin{document}
\maketitle

\section{Problem Statement}
Using all digits 1 through 9 exactly once, how many distinct sets of primes can be formed by concatenating subsets of these digits?

\section{Solution}

\subsection{Approach}
We partition the set $D = \{1, 2, \ldots, 9\}$ into non-empty subsets $S_1, S_2, \ldots, S_k$ such that:
\begin{enumerate}
\item Each subset $S_i$ forms a prime number $p_i$ under some permutation of its digits.
\item $p_1 < p_2 < \cdots < p_k$ (to count sets, not sequences).
\item $\bigcup_{i=1}^k S_i = D$ and the subsets are disjoint.
\end{enumerate}

\subsection{Algorithm}
We use depth-first search with a bitmask representing remaining available digits.

At each recursive step:
\begin{itemize}
\item Enumerate all non-empty subsets of the remaining digits.
\item For each subset, generate all permutations and check if the resulting number is prime.
\item If prime and greater than the last number chosen, recurse with updated bitmask.
\item When the bitmask is empty (all digits used), increment the count.
\end{itemize}

\subsection{Pruning}
\begin{itemize}
\item Numbers whose digit sum is divisible by 3 are divisible by 3 (and thus not prime, unless equal to 3).
\item Even numbers $> 2$ are excluded.
\item Numbers ending in 5 and $> 5$ are excluded.
\end{itemize}

\subsection{Result}
\[
\boxed{44680}
\]

\section{Complexity}
The search space is bounded by the number of ordered partitions of a 9-element set, times the permutations within each part. With aggressive pruning (divisibility by 2, 3, 5), the actual computation is very fast (under 1 second).

\section{Answer}

\[
\boxed{44680}
\]

\end{document}
